Determinar la función a la que converge la serie de Fourier de cada f en el intervalo indicado.
No es necesario calcular los coeficientes. Describir también los valores de la suma en los puntos de salto y empalme. La serie usa el período dado por la longitud del intervalo.

\begin{enumerate}[label=(\alph*),ref=\alph*,leftmargin=*]
\item \(f(x)=x\), \(-\pi<x<\pi\).
\item \(f(x)=|x|\), \(-\pi<x<\pi\).
\item \(f=1\) en(\ensuremath{-}2,0) y \(f=x\) en(0,2).
\item \(f=0\) en(\ensuremath{-}\ensuremath{\pi},0) y \(f=x^2\) en(0,\ensuremath{\pi}).
\item \(f(x)=x^2\), \(-1<x<1\).
\item \(f=x+\pi\) en(\ensuremath{-}\ensuremath{\pi},0) y \(f=x\) en(0,\ensuremath{\pi}).
\item \(f(x)=e^x\), \(-\pi<x<\pi\).
\item \[f(x)=\begin{cases}0,&-\pi<x<-\pi/2,\\-1,&-\pi/2<x<0,\\1,&0<x<\pi/2,\\0,&\pi/2<x<\pi.\end{cases}\]
\end{enumerate}
